\subsection{Full System Specification}

This section presents a consolidated textual specification of the Steam Store Management system. It synthesizes all entities, attributes, relationships, and business rules previously identified in Sections 1.1, 1.2, and 1.3 into a single description. The purpose of this specification is to give a complete picture of the proposed system so that the EERD in Section 2 and the relational schema in Section 3 can be derived and validated against it.

The proposed system is a digital distribution platform inspired by Steam. It must store the entire information of users, games, downloadable content, promotions, transactions, ownership records, reviews, achievements, friendships, and shopping carts. The system supports two main types of users: \textbf{players}, who consume digital content, and \textbf{developers/publishers}, who provide and manage digital products. Both share a common base entity \texttt{USER} that records account-level information such as a unique user identifier, a unique username, a unique email address, a hashed password, and a join date. The specialization between \texttt{PLAYER} and \texttt{DEVELOPER} is total and disjoint: every user must be either a player or a developer, but not both.

A \texttt{PLAYER} extends the base user with player-specific attributes such as a wallet balance and a personal profile bio. The wallet balance must always be greater than or equal to zero and is reduced when the player completes a purchase. A player may own multiple games through the \texttt{OWNS\_LIBRARY} relationship and multiple downloadable contents through the \texttt{OWNS\_DLC} relationship. Players may also become friends with other players through the symmetric \texttt{FRIENDS} relationship, in which the smaller user identifier must always be stored before the larger to avoid duplicate friendship records. Each player may unlock achievements that belong to specific games, and the unlock event is recorded together with the achievement, the player, and the time of unlocking through the \texttt{PLAYER\_ACHIEVEMENT} relationship.

A \texttt{DEVELOPER} extends the base user with a legal name and an optional bank account used for revenue settlement. Each developer may register one or more contact phone numbers, which are stored in the multi-valued attribute table \texttt{CONTACT\_PHONE}. A developer publishes one or more \texttt{GAME} products. Each game is identified by a unique game identifier and includes attributes such as title, base price, release date, system requirements (graphics, operating system, processor, memory), and a derived attribute \texttt{effective\_price} that reflects the current selling price after applying the best valid promotion. A game may belong to one or more genres through the \texttt{GAME\_GENRE} relationship and may be linked to one or more promotions through the \texttt{GAME\_PROMOTION} relationship.

The system also manages \textbf{downloadable content} (\texttt{DLC}) as a weak entity that depends on a parent \texttt{GAME}. Each DLC is identified by a composite key consisting of the parent game identifier and a local DLC identifier. A DLC has its own title, price, and release date. DLCs may be purchased separately from the base game and are tracked independently in the \texttt{OWNS\_DLC} relationship.

\textbf{Promotions} are first-class entities identified by a unique promotion identifier and include attributes such as discount percentage, start date, and end date. The end date must always be greater than or equal to the start date. A promotion may apply to multiple games and a game may receive multiple promotions over time, but at any given moment only the best valid promotion is reflected in the \texttt{effective\_price} of a game.

\textbf{Commercial activity} is divided into two stages: temporary intent and confirmed transaction. The temporary intent is represented by the \texttt{CART} entity, in which a player places games and DLCs that they intend to purchase later. Cart contents are stored through the \texttt{QUEUE} relationship and may be removed without creating a permanent record. The confirmed transaction is represented by the \texttt{TRANSACTION} entity, which is identified by a transaction identifier and includes attributes such as the buyer player, the transaction date, the payment method, and the total amount. Games purchased in a transaction are recorded in the \texttt{CONTAINS} relationship together with the derived attribute \texttt{price\_at\_purchase}, which preserves the exact price paid at the time of purchase regardless of any later price change. DLC purchases inside a transaction are recorded in the \texttt{BUY\_DLC} relationship.

\textbf{Community interaction} is supported through the \texttt{REVIEW} entity. A player may write at most one review for the same game, and only for games that the player already owns. Each review records a binary rating that takes only one of two values: \texttt{Recommended} or \texttt{Not Recommended}, together with the review text and the review date. Reviews are connected to both the player and the game through the \texttt{REVIEW} entity itself.

\textbf{Achievements} are managed through the \texttt{ACHIEVEMENT} weak entity, identified by the parent game and a local achievement identifier. Each achievement has its own name and description. Achievements unlocked by players are tracked through the \texttt{PLAYER\_ACHIEVEMENT} relationship together with the unlock time.

The system is also subject to several \textbf{semantic constraints} that cannot be fully captured in the EERD and must be enforced at the database level through CHECK constraints, triggers, or stored procedure validations. The most important constraints include: the \texttt{wallet\_balance} of a player must remain non-negative; the \texttt{Rating} of a review must be either \texttt{Recommended} or \texttt{Not Recommended}; a player may submit at most one review per game; a player may write a review only for games they own; promotion dates must be valid (\texttt{end\_date} $\geq$ \texttt{start\_date}); transaction totals must equal the sum of \texttt{price\_at\_purchase} of all contained items; and the derived attributes \texttt{effective\_price} and \texttt{price\_at\_purchase} must be automatically computed and maintained whenever related data changes.

Overall, the specification above defines a digital distribution platform that integrates account management, product publication, promotion-aware purchasing, persistent ownership, community-driven reviews, and achievement-based engagement. This consolidated description serves as the conceptual basis for the EERD presented in Section 2 and the relational mapping presented in Section 3.
\newpage
